% ==============================================================================
% CAPITOLO 9: CAMPIONE STATISTICO, TEORIA DELLA STIMA PUNTUALE E MASSIMA VEROSIMIGLIANZA
% ==============================================================================
\chapter{Campione Statistico, Teoria della Stima Puntuale e Massima Verosimiglianza}
\label{cap:campione_statistico_e_stimatori}

\begin{abstract}
Il presente capitolo segna il passaggio concettuale dall'analisi probabilistica diretta (deduttiva: dal modello teorico ai dati attesi) all'inferenza statistica (induttiva: dai dati empirici osservati alla stima dei parametri incogniti del modello). Viene formalizzata la nozione di campione casuale i.i.d. e stabilita la netta distinzione tra statistica, stimatore (variabile aleatoria) e stima (valore numerico realizzato). Si definiscono e si dimostrano i criteri di ottimalità degli stimatori: non distorsione (correttezza), scomposizione dell'Errore Quadratico Medio (MSE) nel trade-off tra distorsione e varianza, consistenza asintotica ed efficienza UMVUE mediante il Limite di Cramér-Rao (CRLB) e l'Informazione di Fisher. Si sviluppa in dettaglio il Metodo della Massima Verosimiglianza (MLE), dimostrandone le proprietà di invarianza e normalità asintotica. Infine, si espone il metodo Monte Carlo per l'integrazione stocastica con relativo codice R.
\end{abstract}

% ==============================================================================
% SEZIONE 9.1: IL PARADIGMA DELL'INFERENZA STATISTICA
% ==============================================================================
\section{Il Paradigma dell'Inferenza: Campione Casuale e Stimatori}
\label{sec:paradigma_inferenza}

In tutte le scienze sperimentali e nell'ingegneria, il processo generativo che governa un fenomeno reale (il tempo di vita di una turbina, il tasso di difettosità di una linea di montaggio, il rendimento energetico di un impianto) è descrivibile mediante una legge di probabilità $f(x; \theta)$ la cui forma funzionale è nota a meno di uno o più parametri incogniti $\theta \in \Theta \subseteq \mathbb{R}^k$.

\begin{definizione}{Campione Casuale I.I.D.}{def_campione_casuale}
Sia $X$ una variabile aleatoria con densità o massa $f(x; \theta)$ dipendente dal parametro $\theta \in \Theta$.
Un \textbf{campione casuale di dimensione $n$} è una $n$-upla di variabili aleatorie indipendenti e identicamente distribuite:
\begin{equation}
    (X_1, X_2, \dots, X_n) \stackrel{\text{i.i.d.}}{\sim} f(x; \theta)
\end{equation}
La realizzazione numerica osservata a posteriori sui dati empirici è denotata con $(x_1, x_2, \dots, x_n) \in \mathbb{R}^n$.
\end{definizione}

\begin{definizione}{Statistica, Stimatore e Stima}{def_stimatore}
\begin{itemize}
    \item Una \textbf{statistica} $T = g(X_1, \dots, X_n)$ è una qualsiasi funzione a valori reali del campione che non contiene parametri incogniti.
    \item Uno \textbf{stimatore puntuale} $\widehat{\theta} = T(X_1, \dots, X_n)$ è una statistica concepita per approssimare il vero valore del parametro $\theta$. Prima dell'estrazione del campione, $\widehat{\theta}$ è una \textbf{variabile aleatoria} dotata di una propria distribuzione campionaria (detta \emph{sampling distribution}).
    \item Una \textbf{stima puntuale} $\widehat{\theta}_{\text{oss}} = T(x_1, \dots, x_n)$ è il singolo numero reale ottenuto sostituendo i dati numerici realizzati nello stimatore.
\end{itemize}
\end{definizione}

% ==============================================================================
% SEZIONE 9.2: CRITERI DI OTTIMALITÀ DEGLI STIMATORI
% ==============================================================================
\section{Criteri di Ottimalità: Correttezza, MSE e Consistenza}
\label{sec:ottimalita_stimatori}

Dato che per uno stesso parametro $\theta$ è possibile costruire molteplici stimatori concorrenti, è fondamentale stabilire criteri oggettivi di qualità.

\subsection{Correttezza (Non Distorsione / Unbiasedness)}

\begin{definizione}{Distorsione e Stimatore Corretto}{def_correttezza}
La \textbf{distorsione} (\emph{Bias}) di uno stimatore $\widehat{\theta}$ è la deviazione sistematica del suo valore atteso dal valore vero del parametro:
\begin{equation}
    \operatorname{Bias}(\widehat{\theta}) \coloneqq \mathbb{E}[\widehat{\theta}] - \theta
\end{equation}
Uno stimatore si dice \textbf{corretto (non distorto)} se il suo valore atteso coincide esattamente con il parametro cercato per ogni possibile $\theta \in \Theta$:
\begin{equation}
    \mathbb{E}[\widehat{\theta}] = \theta \iff \operatorname{Bias}(\widehat{\theta}) = 0
\end{equation}
Si dice \textbf{asintoticamente corretto} se $\lim_{n \to \infty} \mathbb{E}[\widehat{\theta}_n] = \theta$.
\end{definizione}

\subsection{Errore Quadratico Medio e Scomposizione Bias-Varianza}

Non sempre uno stimatore non distorto è preferibile: se esso possiede una varianza enorme, le sue singole stime oscilleranno violentemente lontano dal bersaglio. La misura globale di accuratezza è l'\textbf{Errore Quadratico Medio}.

\begin{figure}[htbp]
    \centering
    \begin{tikzpicture}[scale=0.8, >=Stealth]
    % Bersaglio 1: Basso Bias, Bassa Varianza (Ideale)
    \begin{scope}[shift={(0,0)}]
        \draw[thick, navyblue] (0,0) circle (1.5cm);
        \draw[thin, gray!70] (0,0) circle (1.0cm);
        \draw[thin, gray!70] (0,0) circle (0.5cm);
        \fill[crimsonred] (0,0) circle (2pt) node[below=2pt, font=\tiny\bfseries] {$\theta$};

        \foreach \pt in {(0.1,0.1), (-0.1,-0.15), (0.15,-0.08), (-0.05,0.12), (0.02, -0.05)} {
            \fill[forestgreen] \pt circle (1.8pt);
        }
        \node[below, font=\scriptsize\bfseries, text=forestgreen] at (0,-1.7) {Bias $=0$, Var Bassa};
    \end{scope}

    % Bersaglio 2: Alto Bias, Bassa Varianza
    \begin{scope}[shift={(4.5,0)}]
        \draw[thick, navyblue] (0,0) circle (1.5cm);
        \draw[thin, gray!70] (0,0) circle (1.0cm);
        \draw[thin, gray!70] (0,0) circle (0.5cm);
        \fill[crimsonred] (0,0) circle (2pt) node[below=2pt, font=\tiny\bfseries] {$\theta$};

        \foreach \pt in {(0.8,0.8), (0.9,0.7), (0.75,0.85), (0.85,0.9), (0.7, 0.75)} {
            \fill[warmamber!90!black] \pt circle (1.8pt);
        }
        \node[below, font=\scriptsize\bfseries, text=warmamber!90!black] at (0,-1.7) {Bias Alto, Var Bassa};
    \end{scope}

    % Bersaglio 3: Basso Bias, Alta Varianza
    \begin{scope}[shift={(9.0,0)}]
        \draw[thick, navyblue] (0,0) circle (1.5cm);
        \draw[thin, gray!70] (0,0) circle (1.0cm);
        \draw[thin, gray!70] (0,0) circle (0.5cm);
        \fill[crimsonred] (0,0) circle (2pt) node[below=2pt, font=\tiny\bfseries] {$\theta$};

        \foreach \pt in {(-0.8,0.6), (0.7,0.8), (-0.5,-0.9), (0.8,-0.5), (-0.1, 0.2)} {
            \fill[coolblue] \pt circle (1.8pt);
        }
        \node[below, font=\scriptsize\bfseries, text=coolblue] at (0,-1.7) {Bias $=0$, Var Alta};
    \end{scope}

    % Bersaglio 4: Alto Bias, Alta Varianza (Pessimo)
    \begin{scope}[shift={(13.5,0)}]
        \draw[thick, navyblue] (0,0) circle (1.5cm);
        \draw[thin, gray!70] (0,0) circle (1.0cm);
        \draw[thin, gray!70] (0,0) circle (0.5cm);
        \fill[crimsonred] (0,0) circle (2pt) node[below=2pt, font=\tiny\bfseries] {$\theta$};

        \foreach \pt in {(0.6,0.5), (1.2,1.1), (0.3,1.3), (1.1,0.4), (0.8, 1.0)} {
            \fill[crimsonred] \pt circle (1.8pt);
        }
        \node[below, font=\scriptsize\bfseries, text=crimsonred] at (0,-1.7) {Bias Alto, Var Alta};
    \end{scope}
\end{tikzpicture}

    \caption{Analogia del tiro al bersaglio per il trade-off Distorsione-Varianza (Bias-Variance Tradeoff): (a) Basso Bias e Bassa Varianza (ottimo); (b) Alto Bias e Bassa Varianza (errore sistematico preciso); (c) Basso Bias e Alta Varianza (centrato in media ma disperso); (d) Alto Bias e Alta Varianza (pessimo).}
    \label{fig:bias_variance_tradeoff}
\end{figure}

\begin{teorema}{Decomposizione dell'Errore Quadratico Medio (MSE)}{mse_tradeoff}
L'\textbf{Errore Quadratico Medio} (Mean Squared Error, MSE) di uno stimatore $\widehat{\theta}$ misura la deviazione quadratica totale attesa dal parametro vero $\theta$:
\begin{equation}
    \operatorname{MSE}(\widehat{\theta}) \coloneqq \mathbb{E}\left[(\widehat{\theta} - \theta)^2\right] = \operatorname{Var}(\widehat{\theta}) + \left[\operatorname{Bias}(\widehat{\theta})\right]^2
\end{equation}
Se lo stimatore è corretto ($\operatorname{Bias}=0$), l'MSE coincide con la varianza: $\operatorname{MSE}(\widehat{\theta}) = \operatorname{Var}(\widehat{\theta})$.
\end{teorema}

\begin{dimostrazione}
Aggiungendo e sottraendo il valore atteso $\mathbb{E}[\widehat{\theta}]$ all'interno dell'operatore:
\begin{align}
    \operatorname{MSE}(\widehat{\theta}) &= \mathbb{E}\left[\left( (\widehat{\theta} - \mathbb{E}[\widehat{\theta}]) + (\mathbb{E}[\widehat{\theta}] - \theta) \right)^2\right] \\
    &= \mathbb{E}\left[(\widehat{\theta} - \mathbb{E}[\widehat{\theta}])^2\right] + 2(\mathbb{E}[\widehat{\theta}] - \theta)\underbrace{\mathbb{E}\left[\widehat{\theta} - \mathbb{E}[\widehat{\theta}]\right]}_{=0} + (\mathbb{E}[\widehat{\theta}] - \theta)^2 \\
    &= \operatorname{Var}(\widehat{\theta}) + \left[\operatorname{Bias}(\widehat{\theta})\right]^2 \qquad \blacksquare
\end{align}
\end{dimostrazione}

\subsection{Consistenza Asintotica}

\begin{definizione}{Stimatore Consistente}{def_consistenza}
Una successione di stimatori $\{\widehat{\theta}_n\}_{n \ge 1}$ si dice \textbf{consistente in probabilità} (o debolmente consistente) per $\theta$ se all'aumentare della numerosità campionaria $n \to \infty$ converge in probabilità al parametro vero:
\begin{equation}
    \forall \varepsilon > 0, \quad \lim_{n \to \infty} \mathbb{P}\left(|\widehat{\theta}_n - \theta| \ge \varepsilon\right) = 0 \iff \widehat{\theta}_n \xrightarrow{\mathbb{P}} \theta
\end{equation}
\textbf{Condizione Sufficiente di Consistenza (tramite Chebyshev):} Se lo stimatore è asintoticamente non distorto e la sua varianza tende a zero per $n \to \infty$, allora è consistente:
\begin{equation}
    \lim_{n \to \infty} \operatorname{Bias}(\widehat{\theta}_n) = 0 \quad \text{e} \quad \lim_{n \to \infty} \operatorname{Var}(\widehat{\theta}_n) = 0 \implies \lim_{n \to \infty} \operatorname{MSE}(\widehat{\theta}_n) = 0 \implies \widehat{\theta}_n \xrightarrow{\mathbb{P}} \theta
\end{equation}
\end{definizione}

% ==============================================================================
% SEZIONE 9.3: LIMITE DI CRAMÉR-RAO E INFORMAZIONE DI FISHER
% ==============================================================================
\section{Informazione di Fisher e Limite di Cramér-Rao (CRLB)}
\label{sec:fisher_cramer_rao}

Esiste un limite fisico insormontabile alla precisione con cui è possibile stimare un parametro da un campione di taglia $n$? La risposta è fornita dalla teoria dell'Informazione di Fisher.

\begin{definizione}{Informazione di Fisher}{def_informazione_fisher}
Sia $f(x; \theta)$ la densità della popolazione. Sotto condizioni di regolarità analitica, l'\textbf{Informazione di Fisher} contenuta in una singola osservazione è definita da:
\begin{equation}
    I(\theta) \coloneqq \mathbb{E}\left[ \left( \frac{\partial}{\partial \theta} \ln f(X; \theta) \right)^2 \right] = -\mathbb{E}\left[ \frac{\partial^2}{\partial \theta^2} \ln f(X; \theta) \right]
\end{equation}
Per un campione i.i.d. di dimensione $n$, l'informazione totale è additiva: $I_n(\theta) = n I(\theta)$.
\end{definizione}

\paragraph{Significato Geometrico dell'Informazione di Fisher:}
La derivata seconda $-\frac{\partial^2}{\partial \theta^2} \ln f(x; \theta)$ misura la **curvatura (o concavità)** della log-verosimiglianza. Se la curva è molto piccata e incurvata verso il basso attorno al massimo ($I(\theta)$ grande), piccole variazioni di $\theta$ alterano drasticamente la densità, consentendo una stima estremamente precisa; se la curva è piatta ($I(\theta)$ piccolo), vi è alta incertezza.

\begin{teorema}{Disuguaglianza e Limite Inferiore di Cramér-Rao (CRLB)}{crlb}
Sia $\widehat{\theta}$ uno stimatore corretto ($\mathbb{E}[\widehat{\theta}] = \theta$) per il parametro scalare $\theta$. Sotto condizioni di regolarità, la varianza dello stimatore è limitata inferiormente dal reciproco dell'informazione di Fisher totale:
\begin{equation}
    \operatorname{Var}(\widehat{\theta}) \ge \frac{1}{I_n(\theta)} = \frac{1}{n I(\theta)} = \text{CRLB}
\end{equation}
Uno stimatore corretto che raggiunge esattamente l'uguaglianza $\operatorname{Var}(\widehat{\theta}) = \text{CRLB}$ si dice \textbf{efficiente} o a varianza minima uniformemente (UMVUE).
\end{teorema}

% ==============================================================================
% SEZIONE 9.4: METODO DELLA MASSIMA VEROSIMIGLIANZA (MLE)
% ==============================================================================
\section{Il Metodo della Massima Verosimiglianza (MLE)}
\label{sec:mle_metodo}

Il Metodo della Massima Verosimiglianza (\emph{Maximum Likelihood Estimation - MLE}), introdotto da Ronald Fisher nel 1922, è la tecnica di stima parametrica più autorevole ed elegante dell'inferenza statistica.

\begin{figure}[htbp]
    \centering
    \begin{tikzpicture}[scale=0.85, >=Stealth]
    \draw[->, thick, navyblue] (-0.5,0) -- (6.5,0) node[right, font=\footnotesize] {$\theta$};
    \draw[->, thick, navyblue] (0,-0.3) -- (0,3.8) node[above, font=\footnotesize] {$\ell(\theta) = \ln L(\theta)$};

    % Curva log-verosimiglianza (concava)
    \draw[domain=0.8:5.5, smooth, variable=\x, ultra thick, navyblue]
        plot ({\x}, {3.2 - 0.7*(\x-3.2)*(\x-3.2)});

    % Massimo a theta = 3.2
    \draw[dashed, crimsonred, thick] (3.2,0) -- (3.2,3.2);
    \fill[crimsonred] (3.2,3.2) circle (2.5pt);
    \node[below, font=\small\bfseries, text=crimsonred] at (3.2,0) {$\widehat{\theta}_{\text{MLE}}$};

    % Tangente orizzontale
    \draw[thick, forestgreen] (2.0,3.2) -- (4.4,3.2);
    \node[above, font=\scriptsize\bfseries, text=forestgreen] at (3.2,3.3) {Score $= \dfrac{d\ell}{d\theta} = 0$};

    \node[font=\footnotesize\itshape, text=navyblue, align=center] at (3.2, -0.9) {
        Curvatura nel massimo $= -J_n(\widehat{\theta}) = \left.\dfrac{d^2\ell}{d\theta^2}\right|_{\widehat{\theta}} < 0$ (\textbf{Informazione di Fisher}).
    };
\end{tikzpicture}

    \caption{Curva della log-verosimiglianza $\ell(\theta)$ in funzione del parametro $\theta$. Lo stimatore MLE $\widehat{\theta}_{\text{MLE}}$ individua il punto stazionario di massimo assoluto in cui la funzione score si annulla: $S(\theta) = \ell'(\theta) = 0$, con derivata seconda strettamente negativa $\ell''(\theta) < 0$.}
    \label{fig:mle_curva}
    \end{figure}

\begin{definizione}{Funzione di Verosimiglianza e Log-Verosimiglianza}{def_verosimiglianza}
Dato il campione osservato $\mathbf{x} = (x_1, \dots, x_n)$, la \textbf{funzione di verosimiglianza} $L(\theta)$ è la densità congiunta interpretata come funzione del parametro $\theta$:
\begin{equation}
    L(\theta) \coloneqq f(x_1, \dots, x_n; \theta) = \prod_{i=1}^n f(x_i; \theta)
\end{equation}
Poiché il logaritmo naturale è una trasformazione strettamente monotona crescente, la massimizzazione di $L(\theta)$ è matematicamente equivalente alla massimizzazione della \textbf{log-verosimiglianza}:
\begin{equation}
    \ell(\theta) \coloneqq \ln L(\theta) = \sum_{i=1}^n \ln f(x_i; \theta)
\end{equation}
\end{definizione}

\begin{metodo}[Procedura Operativa in 4 Fasi per il Calcolo dello Stimatore MLE]
\begin{enumerate}
    \item \textbf{Scrittura della Verosimiglianza:} $L(\theta) = \prod_{i=1}^n f(x_i; \theta)$.
    \item \textbf{Passaggio al Logaritmo:} $\ell(\theta) = \sum_{i=1}^n \ln f(x_i; \theta)$.
    \item \textbf{Calcolo della Score Function ed Uguaglianza a Zero:}
    \begin{equation}
        S(\theta) = \frac{d\ell(\theta)}{d\theta} = 0 \implies \text{si risolve l'equazione per ricavare } \widehat{\theta}
    \end{equation}
    \item \textbf{Verifica della Concavità (Test dell'Hessiano):}
    \begin{equation}
        \frac{d^2\ell(\theta)}{d\theta^2} \Bigg|_{\theta = \widehat{\theta}} < 0 \implies \text{garantisce che il punto stazionario sia un massimo relativo e globale.}
    \end{equation}
\end{enumerate}
\end{metodo}

\begin{teorema}{Proprietà Notevoli dello Stimatore MLE}{prop_mle}
\begin{enumerate}
    \item \textbf{Proprietà di Invarianza Funzionale (Teorema di Zehna):} Se $\widehat{\theta}_{\text{MLE}}$ è lo stimatore di massima verosimiglianza per $\theta$, allora per qualsiasi funzione $g(\theta)$, lo stimatore MLE di $g(\theta)$ è semplicemente $g(\widehat{\theta}_{\text{MLE}})$.
    \item \textbf{Normalità Asintotica ed Efficienza:} All'aumentare di $n \to \infty$, lo stimatore MLE converge in legge ad una normale con varianza pari al limite di Cramér-Rao:
    \begin{equation}
        \sqrt{n}(\widehat{\theta}_{\text{MLE}} - \theta) \xrightarrow{d} \mathcal{N}\left(0, \; \frac{1}{I(\theta)}\right)
    \end{equation}
\end{enumerate}
\end{teorema}

% ==============================================================================
% SEZIONE 9.5: METODO MONTE CARLO IN R
% ==============================================================================
\section{Il Metodo Monte Carlo per la Stima e l'Integrazione Numerica}
\label{sec:monte_carlo_teoria}

Quando un integrale multivariato o il valore atteso di una grandezza complessa $I = \mathbb{E}[g(X)] = \int g(x) f(x) dx$ non ammette una primitiva in forma chiusa, il metodo Monte Carlo sfrutta la Legge dei Grandi Numeri per stimare $I$ mediante campionamento pseudo-casuale al calcolatore.

\begin{figure}[htbp]
    \centering
    \begin{tikzpicture}[scale=0.85, >=Stealth]
    % Quadrato [0,1]^2
    \draw[thick, navyblue] (0,0) rectangle (4,4);
    \node[below left, font=\tiny] at (0,0) {$(0,0)$};
    \node[below, font=\tiny] at (4,0) {$1$};
    \node[left, font=\tiny] at (0,4) {$1$};

    % Quarto di cerchio unitario
    \draw[ultra thick, crimsonred] (4,0) arc (0:90:4);
    \fill[coolblue!15] (0,0) -- (4,0) arc (0:90:4) -- cycle;

    % Punti campionati
    \foreach \pt in {(0.5,0.8), (1.2,1.5), (2.0,1.8), (2.5,2.2), (0.8,3.1), (1.5,2.8), (3.0,1.2), (0.3,2.0), (1.8,0.5), (2.8,0.9)} {
        \fill[forestgreen] \pt circle (2pt);
    }
    \foreach \pt in {(3.2,3.1), (3.6,2.5), (2.8,3.5), (3.8,1.8), (1.8,3.8), (3.5,3.7)} {
        \fill[crimsonred] \pt circle (2pt);
    }

    \node[font=\scriptsize\bfseries, text=navyblue, align=center] at (6.5,2.0) {
        $\dfrac{N_{\text{dentro}}}{N_{\text{tot}}} \approx \dfrac{\text{Area}(\text{Cerchio})}{\text{Area}(\text{Quadrato})} = \dfrac{\pi}{4}$ \\[3mm]
        $\widehat{\pi} = 4 \cdot \dfrac{N_{\text{dentro}}}{N_{\text{tot}}}$
    };
\end{tikzpicture}

    \caption{Principio del metodo Monte Carlo: stima geometrica di $\pi$ tramite campionamento uniforme di punti $(X_i, Y_i) \sim \operatorname{Unif}([-1,1]^2)$. La frazione di punti che cade all'interno del cerchio unitario $X^2 + Y^2 \le 1$ converge a $\frac{\text{Area Cerchio}}{\text{Area Quadrato}} = \frac{\pi}{4}$.}
    \label{fig:monte_carlo_cerchio}
\end{figure}

\begin{metodo}[Integrazione Monte Carlo in R]
Per approssimare l'integrale $I = \int_a^b h(x) dx$:
\begin{enumerate}
    \item Si riscrive l'integrale come $I = (b - a) \int_a^b h(x) \frac{1}{b-a} dx = (b - a) \mathbb{E}[h(X)]$ con $X \sim \operatorname{Unif}(a, b)$.
    \item Si generano $N$ numeri pseudo-casuali $x_1, \dots, x_N \sim \operatorname{Unif}(a, b)$ e si calcola la media:
    \begin{equation}
        \widehat{I}_N = (b - a) \cdot \frac{1}{N}\sum_{i=1}^N h(x_i)
    \end{equation}
\end{enumerate}
L'errore standard della stima decade come $\mathcal{O}(1/\sqrt{N})$, indipendentemente dalla dimensione dello spazio di integrazione.
\end{metodo}
